\documentclass[conference]{IEEEtran}
% Packages used
\usepackage{amsmath,graphicx}
\usepackage{color}
\usepackage{graphicx}
\usepackage{epstopdf}
\usepackage{amsmath}
\usepackage{amssymb}
\usepackage{mathrsfs}
%\usepackage{dutchcal}
\usepackage{hyperref}
%\usepackage{amsthm}
\usepackage{tikz}
%\usepackage[bbgreekl]{mathbbol}
%\DeclareSymbolFontAlphabet{\mathbbm}{bbold}
\usepackage{bm}
\usepackage[english]{babel}
\usepackage{cite}
\usepackage{rotfloat}
\usepackage{mathtools}
\usepackage[font=normalsize,labelfont=bf]{caption}
%\usepackage{subcaption}
\usepackage{empheq}
\newcommand*\widefbox[1]{\fbox{\hspace{2em}#1\hspace{2em}}}

%\usepackage{breqn}
\usepackage{amsmath}
\usepackage{makecell}
\usepackage{algorithm,algorithmic}
%\usepackage[ruled, vlined, linesnumbered]{algorithm2e}
\usepackage{multirow}
\usepackage{subfigure}
\usepackage{booktabs}
\usepackage{colortbl}
\usepackage{multirow}% http://ctan.org/pkg/multirow
\usepackage{hhline}% http://ctan.org/pkg/hhline
\usepackage{stfloats}% <-- added
\usepackage{multicol}
\usepackage{bbm}
\usepackage{cases}
\graphicspath{ {Figures/} }
\setlength{\textfloatsep}{0.1cm}
\DeclarePairedDelimiter\ceil{\lceil}{\rceil}
\DeclarePairedDelimiter\floor{\lfloor}{\rfloor}
\setlength{\abovecaptionskip}{0.1cm} % Chosen fairly arbitrarily
\newcommand{\eqtext}[1]{\ensuremath{\stackrel{\text{#1}}{=}}}

% transpose, hermitian transpose
\newcommand{\T}{{\scriptscriptstyle\mathsf{T}}}
\renewcommand{\H}{{\scriptscriptstyle\mathsf{H}}}

\newsavebox{\foobox}
\newcommand{\slantbox}[2][.3]
{%
	\mbox
	{%
		\sbox{\foobox}{#2}%
		\hskip\wd\foobox
		\pdfsave
		\pdfsetmatrix{1 0 #1 1}%
		\llap{\usebox{\foobox}}%
		\pdfrestore
	}%
}
\newcommand{\setA}{\mathcal{A}}

\definecolor{kugray5}{RGB}{224,224,224}

\usepackage[normalem]{ulem}
\newcommand\rsout{\bgroup\markoverwith
	{\textcolor{red}{\rule[0.5ex]{2pt}{0.8pt}}}\ULon}

%% ALLOW TO DIVIDE AN ALGORITHM INTO TWO PAGES
\makeatletter
\newenvironment{breakablealgorithm}
{% \begin{breakablealgorithm}
		\begin{center}
			\refstepcounter{algorithm}% New algorithm
			\hrule height.8pt depth0pt \kern2pt% \@fs@pre for \@fs@ruled
			\renewcommand{\caption}[2][\relax]{% Make a new \caption
				{\raggedright\textbf{\ALG@name~\thealgorithm} ##2\par}%
				\ifx\relax##1\relax % #1 is \relax
				\addcontentsline{loa}{algorithm}{\protect\numberline{\thealgorithm}##2}%
				\else % #1 is not \relax
				\addcontentsline{loa}{algorithm}{\protect\numberline{\thealgorithm}##1}%
				\fi
				\kern2pt\hrule\kern2pt
			}
		}{% \end{breakablealgorithm}
		\kern2pt\hrule\relax% \@fs@post for \@fs@ruled
	\end{center}
}
\makeatother

% Color definition
\newcommand{\blue}[1]{\textcolor{blue}{#1}}
\newcommand{\red}[1]{\textcolor{red}{#1}}
\newcommand{\green}[1]{\textcolor{green}{#1}}
\newcommand{\orange}[1]{\textcolor{orange}{#1}}

% vector
\newcommand{\myvec}[1]{\ensuremath{\begin{pmatrix}#1\end{pmatrix}}}

% Algorithmic modifications
\makeatletter
\newcommand{\ALOOP}[1]{\ALC@it\algorithmicloop\ #1%
	\begin{ALC@loop}}
	\newcommand{\ENDALOOP}{\end{ALC@loop}\ALC@it\algorithmicendloop}
\renewcommand{\algorithmicrequire}{\textbf{Input:}}
\renewcommand{\algorithmicensure}{\textbf{Output:}}
\newcommand{\algorithmicbreak}{\textbf{break}}
\newcommand{\BREAK}{\STATE \algorithmicbreak}
\makeatother

% Folder of all figures
\usepackage{etoolbox}
\let\mybibitem\bibitem
\renewcommand{\bibitem}[1]{%
	\ifstrequal{#1}{nature}
	{\color{blue}\mybibitem{#1}}
	{\color{black}\mybibitem{#1}}%
}

\graphicspath{ {Figures/} }

% To define the theorems, definitions, remark, lemma, corollary, proof
\newtheorem{definition}{\textbf{Definition}}
\newtheorem{theorem}{\textbf{Theorem}}
\newtheorem{remark}{\textbf{Remark}}
\newtheorem{lemma}{\textbf{Lemma}}
\newtheorem{corollary}{\textbf{Corollary}}
\newtheorem{proof}{Proof}
\renewcommand\theproof{\unskip}
\newcommand{\epr}{\hfill\(\Box\)}

% for 0.1cm and 0.25cm space
\newcommand{\dis}{\hspace{0.1cm}}
\DeclareCaptionLabelSeparator{periodspace}{.\quad}

% for format of Fig. in IEEE
\renewcommand{\figurename}{Fig.}
% to be able to place 2 figure consecutive at the top of page
\captionsetup{font=footnotesize,labelsep=periodspace,singlelinecheck=false}
\captionsetup[sub]{font=footnotesize,singlelinecheck=true}
\addto\captionsenglish{\renewcommand{\figurename}{Fig.}}
% to be able to palce appendices right after conclusion
\interdisplaylinepenalty=2500 

\newcommand\nbthis{\addtocounter{equation}{1}\tag{\theequation}}

% math notations and notations
\newcommand{\norm}[1]{\left\lVert#1\right\rVert} % ||.||
\newcommand{\fronorm}[1]{\left\lVert#1\right\rVert_{\mathcal{F}}} % ||.||
\newcommand{\normshort}[1]{\lVert#1\rVert} % ||.||
\newcommand{\eq}[1]{\begin{align*}#1\end{align*}} % equation
\newcommand{\eqn}[1]{\begin{align}#1\end{align}} % equation
\newcommand{\nt}[1]{\left(#1\right)} % ()
\newcommand{\nv}[1]{\left[#1\right]} % []
\newcommand{\nn}[1]{\left\{#1\right\}} % {}
\newcommand{\abs}[1]{\left|#1\right|} % ||
\newcommand{\absshort}[1]{|#1|} % ||
\newcommand{\nb}{\numberthis}
\newcommand{\tr}[1]{\mathrm{trace}\left(#1\right)} % ||
\newcommand{\detm}[1]{\mathrm{det}\left(#1\right)} % ||
\newcommand{\trsq}[1]{\mathrm{trace}^2\left(#1\right)} % ||
\newcommand{\diag}[1]{\mathrm{diag}\left\{#1\right\}} % ||
\newcommand{\var}[1]{\mathbb{V}\mathrm{ar}\left(#1\right)} % ||
\newcommand{\ds}{\displaystyle}
% real and imag part
\newcommand{\re}[1]{\mathfrak{R}{\left(#1\right)}}
\newcommand{\im}[1]{\mathfrak{I}{\left(#1\right)}}
\allowdisplaybreaks
% Prob, PDF, CDF
%\newcommand{\prob}{\mathbb{P}} 
\newcommand{\cdf}{\mathbf{\textit{F}}} 
\newcommand{\diagtt}{\mathtt{diag}}
\newcommand{\pdf}{\mathbf{\textit{f}}} 
\newcommand{\mean}[1]{\mathbb{E} \left\{#1\right\}}
\newcommand{\cov}[1]{\mathbb{C} \left\{#1\right\}}
\newcommand{\meanBig}[1]{\mathbb{E} \Big\{#1\Big\}}
\newcommand{\covBig}[1]{\mathbb{C} \Big\{#1\Big\}}

\usepackage{setspace}
% Problem name
\newcommand{\phybrid}{\left(\mathrm{P}_{\mathrm{hybrid}} \right)}
\newcommand{\pnhybrid}{\left(\mathrm{P-n}_{\mathrm{hybrid}} \right)}
\newcommand{\pnbhybrid}{\left(\var{\mathrm{Pn}}_{\mathrm{hybrid}} \right)}
\newcommand{\ppassive}{\left(\mathrm{P}_{\mathrm{passive}} \right)}
\newcommand{\pactive}{\left(\mathrm{P}_{\mathrm{active}} \right)}
\newcommand{\pnpassive}{\left(\mathrm{P-n}_{\mathrm{passive}} \right)}
\newcommand{\pnuniform}{\left(\mathrm{P-n}_{\mathrm{uniform}} \right)}
\newcommand{\pnhybridLB}{\left(\mathrm{P-n}_{\mathrm{fixed, uni.}} \right)}
\newcommand{\pdyn}{\left(\mathrm{P}_{\mathrm{dyn}} \right)}
\newcommand{\pdyna}{\left(\mathrm{P1}_{\mathrm{dyn}} \right)}
\newcommand{\prob}{\left(\mathrm{P} \right)}
\newcommand{\probn}{\left(\mathrm{Pn} \right)}
% Parameters

% Matrix
\newcommand{\mQ}{{\mathbf{Q}}}
\newcommand{\mR}{{\mathbf{R}}}
\newcommand{\mH}{{\mathbf{H}}} 
\newcommand{\mhsorted}{\underline{\mH}}
\newcommand{\mA}{{\mathbf{A}}}
\newcommand{\mW}{{\mathbf{W}}}
\newcommand{\mP}{{\mathbf{P}}}
\newcommand{\mI}{\mathbf{I}}
\newcommand{\mT}{{\mathbf{T}}}
\newcommand{\mB}{{\mathbf{B}}}
\newcommand{\mC}{{\mathbf{C}}}
\newcommand{\mD}{{\mathbf{D}}}
\newcommand{\mX}{{\mathbf{X}}}
\newcommand{\mY}{{\mathbf{Y}}}
\newcommand{\mG}{{\mathbf{G}}}
\newcommand{\mF}{{\mathbf{F}}}
\newcommand{\mU}{{\mathbf{U}}}
\newcommand{\mV}{{\mathbf{V}}}
\newcommand{\mE}{{\mathbf{E}}}
\newcommand{\mHr}{{\mathbf{H}}_r} 
\newcommand{\mHt}{{\mathbf{H}}_t} 
\newcommand{\mIr}{{\mathbf{I}}_{N_r}}
\newcommand{\mZ}{{\mathbf{Z}}}
\newcommand{\mK}{{\mathbf{K}}}
\newcommand{\mJ}{{\mathbf{J}}}
\newcommand{\Na}{N_{\mathrm{A}}}
\def\mathbi#1{\textbf{\em #1}}

% max value
\newcommand{\Xm}{X_{max}} 
\newcommand{\Am}{A_{max}} 
\newcommand{\Bm}{B_{min}} 


%% Set
\newcommand{\setC}{\mathbb{C}} 
\newcommand{\setR}{\mathbb{R}}
%\newcommand{\setA}{\mathcal{A}} 
\newcommand{\setD}{\mathcal{D}} 
\newcommand{\setANt}{\setA^{N_t}} 
\newcommand{\setAtld}{\tilde{\setA}^{N_t}}
\newcommand{\setN}{\mathbb{N}}

%% List
\newcommand{\ltabu}{\mathcal{L}}
\newcommand{\leta}{\mathcal{L}_{\eta}}
\newcommand{\lphi}{\mathcal{L}_{\phi}}

%% Probability
\newcommand{\Pe}{\mathcal{P}}
\newcommand{\Pne}{\var{\mathcal{P}}}

%% Vector
\newcommand{\vxb}{{\mathbf{x}}^{\star}}
\newcommand{\vc}{{\mathbf{c}}}
\newcommand{\vci}{{\mathbf{c}}_{\nn{1}}}
\newcommand{\ve}{{\mathbf{e}}} 
\newcommand{\vs}{{\mathbf{s}}}
\newcommand{\vx}{{\mathbf{x}}}
\newcommand{\vy}{{\mathbf{y}}}
\newcommand{\vr}{{\mathbf{r}}}
\newcommand{\vqn}{{\mathbf{q}}^T_n}
\newcommand{\vpnT}{{\mathbf{p}}^T_n}
%\newcommand{\vpnT}{{\mathbf{r}}^T_n}
\newcommand{\vpd}{{\mathbf{p}}_d}
\newcommand{\vrdc}{\check{{\mathbf{r}}}_d}
\newcommand{\vhdc}{\check{{\mathbf{h}}}_d}
\newcommand{\vhc}{\check{{\mathbf{h}}}}
\newcommand{\vrd}{{\mathbf{r}}_d}
\newcommand{\vrds}{{\mathbf{r}}_{d^{\star}}}
\newcommand{\vpn}{{\mathbf{p}}_n}
\newcommand{\vv}{{\mathbf{v}}}
\newcommand{\vn}{{\mathbf{n}}}
\newcommand{\vu}{{\mathbf{u}}}
\newcommand{\vz}{{\mathbf{z}}} 
\newcommand{\vh}{{\mathbf{h}}} 
\newcommand{\vhd}{{\mathbf{h}}_d}
\newcommand{\vq}{{\mathbf{q}}}
\newcommand{\vb}{{\mathbf{b}}}
\newcommand{\vw}{{\mathbf{w}}}
\newcommand{\va}{{\mathbf{a}}}
\newcommand{\vd}{{\mathbf{d}}}
\newcommand{\vt}{{\mathbf{t}}}
\newcommand{\vg}{{\mathbf{g}}}
\newcommand{\vp}{{\mathbf{p}}}

%% Elements of vectors/matrice
\newcommand{\zbn}{\var{z}_n}
\newcommand{\rni}{r_{n,i}}
\newcommand{\rnj}{r_{n,j}}
\newcommand{\pnj}{p_{n,j}}
\newcommand{\qni}{q_{n,i}}
\newcommand{\qnj}{q_{n,j}}

%% Other notations
\newcommand{\dx}{\Delta{\mathbf{x}}}

\newcommand{\mpeta}[1]{\phi_{#1} (\vx)} % Metric
\newcommand{\mpetai}[1]{\phi_{#1} (\vx_{[i]})} % Metric

\newcommand{\smt}{\sigma_t^2} % Power
\newcommand{\smv}{\sigma_v^2} % Power
\newcommand{\smn}{\sigma_n^2} % Power

\newcommand{\pbar}{\var{\Phi}}
\newcommand{\ptilde}{\tilde{\Phi}}

\newcommand{\bPhi}{\boldsymbol{\Phi}}
\newcommand{\bUpsilon}{\boldsymbol{\Upsilon}}
\newcommand{\bTheta}{\boldsymbol{\Theta}}
\newcommand{\bPsi}{\boldsymbol{\Psi}}
\newcommand{\bSigma}{\boldsymbol{\Sigma}}

\newcommand{\Pa}{P_{\mathrm{a}}} 
\newcommand{\Pamax}{P_{\mathrm{a}}^{\mathrm{max}}}
\newcommand{\Pap}{P_{\mathrm{AP}}}
\newcommand{\amn}{\alpha_{mn}}
\newcommand{\amns}{\alpha_{mn}^{\star}}
\newcommand{\vphi}{\boldsymbol{\varphi}}
\newcommand{\vmu}{\boldsymbol{\mu}}
\newcommand{\vgamma}{\var{\boldsymbol{\gamma}}}
\newcommand{\veta}{\var{\boldsymbol{\eta}}}
\newcommand{\vf}{\mathbf{f}}
\newcommand{\vnu}{\boldsymbol{\nu}}
\newcommand{\vomega}{\boldsymbol{\omega}}
\newcommand{\valpha}{\boldsymbol{\alpha}}
\newcommand{\mSigma}{\boldsymbol{\Sigma}}



% Calligraphic uppercase
\def\cA{{\mathcal{A}}} \def\cB{{\mathcal{B}}} \def\cC{{\mathcal{C}}} \def\cD{{\mathcal{D}}}
\def\cE{{\mathcal{E}}} \def\cF{{\mathcal{F}}} \def\cG{{\mathcal{G}}} \def\cH{{\mathcal{H}}}
\def\cI{{\mathcal{I}}} \def\cJ{{\mathcal{J}}} \def\cK{{\mathcal{K}}} \def\cL{{\mathcal{L}}}
\def\cM{{\mathcal{M}}} \def\cN{{\mathcal{N}}} \def\cO{{\mathcal{O}}} \def\cP{{\mathcal{P}}}
\def\cQ{{\mathcal{Q}}} \def\cR{{\mathcal{R}}} \def\cS{{\mathcal{S}}} \def\cT{{\mathcal{T}}}
\def\cU{{\mathcal{U}}} \def\cV{{\mathcal{V}}} \def\cW{{\mathcal{W}}} \def\cX{{\mathcal{X}}}
\def\cY{{\mathcal{Y}}} \def\cZ{{\mathcal{Z}}} \def\cz{{\mathcal{z}}}

\def\redmark{\textcolor{red}}
\def\argmin{\mathop{\mathrm{argmin}}}
\def\argmax{\mathop{\mathrm{argmax}}}
%\def\diag{\mathop{\mathrm{diag}}}
\def\inf{\mathop{\mathrm{in}}}
\def\outf{\mathop{\mathrm{out}}}
\def\trace{\mathop{\text{tr}}}
%\def\span{\mathop{\mathrm{span}}}
\def\dim{\mathop{\mathrm{dim}}}
\def\Re{\mathop{\mathrm{Re}}}
\def\Im{\mathop{\mathrm{Im}}}
\def\redmark{\textcolor{red}}

\newcommand{\Ei}{{\mathrm{Ei}}}

\def\bDelta{{\pmb{\Delta}}} \def\bdelta{{\pmb{\delta}}}
\def\bSigma{{\pmb{\Sigma}}} \def\bsigma{{\pmb{\sigma}}}
\def\bPhi{{\pmb{\Phi}}} \def\bphi{{\pmb{\phi}}}
\def\bGamma{{\pmb{\Gamma}}} \def\bgamma{{\pmb{\gamma}}}
\def\bOmega{{\pmb{\Omega}}} \def\bomega{{\pmb{\omega}}}
\def\bTheta{{\pmb{\Theta}}} \def\btheta{{\pmb{\theta}}}
\def\bepsilon{{\pmb{\epsilon}}} \def\bPsi{{\pmb{\Psi}}}
\def\obH{\overline{\bH}}\def\obw{\overline{\bw}} \def\obz{\overline{\bz}}\def\bmu{{\pmb{\mu}}}
\def\b0{{\pmb{0}}}\def\bLambda{{\pmb{\Lambda}}} \def\oc{\overline{\bc}}

\newcommand{\Frf}{\mF_\mathrm{RF}}
\newcommand{\Fopt}{\mF_{\text{opt}}}

\newcommand{\Frft}{\tilde{\mF}_\mathrm{RF}}
\newcommand{\Fbb}{\mF_\mathrm{BB}}
\newcommand{\fbb}{\vf_\mathrm{BB}}

\newcommand{\Wrf}{\bW_\mathrm{RF}}
\newcommand{\Wbb}{\bW_\mathrm{BB}}
\newcommand{\Nrf}{N_{\text{RF}}} 
\newcommand{\Nr}{N_\mathrm{r}}
\newcommand{\Nt}{N_\mathrm{t}} 
\newcommand{\Nu}{N_\mathrm{u}}
\newcommand{\Ntn}{N_{\text{t},n}} 
\newcommand{\Ntnt}{\tilde{N}_{\text{t},n}} 
\newcommand{\Ns}{N_\mathrm{s}}
\newcommand{\Nrh}{N_\mathrm{r}^\mathrm{h}}
\newcommand{\Nrv}{N_\mathrm{r}^\mathrm{v}}
\newcommand{\Nth}{N_\mathrm{t}^\mathrm{h}}
\newcommand{\Ntv}{N_\mathrm{t}^\mathrm{v}}

\newcommand{\Nrhn}{N_{\text{r},n}^\mathrm{h}}
\newcommand{\Nrvn}{N_{\text{r},n}^\mathrm{v}}
\newcommand{\Nthn}{N_{\text{t},n}^\mathrm{h}}
\newcommand{\Ntvn}{N_{\text{t},n}^\mathrm{v}}



\newcommand{\Nrfr}{N_{\text{RF,r}}}
\newcommand{\Nrft}{N_{\text{RF,t}}}
\newcommand{\Mnr}{M_{n,\text{r}}}
\newcommand{\Mnt}{M_{n,\text{t}}}
\newcommand{\ar}{\va_\mathrm{r}} 
\newcommand{\at}{\va_\mathrm{t}} 
\newcommand{\arh}{\va_\mathrm{r,h}} 
\newcommand{\arv}{\va_\mathrm{r,v}} 
\newcommand{\ath}{\va_\mathrm{t,h}} 
\newcommand{\atv}{\va_\mathrm{t,v}}
\newcommand{\art}{\tilde{\va}_\mathrm{r}} 
\newcommand{\att}{\tilde{\va}_\mathrm{t}} 
\newcommand{\atht}{\tilde{\va}_\mathrm{t,h}} 
\newcommand{\atvt}{\tilde{\va}_\mathrm{t,v}}

\newcommand{\Afull}{\mathcal{A}_{\text{full}}}
\newcommand{\Asub}{\mathcal{A}_{\text{sub}}}

\newcommand{\setK}{\mathcal{K}}
\newcommand{\setT}{\mathcal{T}}
\newcommand{\setJ}{\mathcal{J}}
\newcommand{\setM}{\mathcal{M}}
\newcommand{\setTheta}{\bm{\Theta_{\mathrm{d}}}}
\newcommand{\Pt}{P_{\mathrm{t}}}
\newcommand{\noise}{\sigma_{\mathrm{n}}^2}
\newcommand{\Iin}{\mathcal{I}_{\mathrm{i}}}
 \newcommand{\Iout}{\mathcal{I}_{\mathrm{o}}}
\IEEEoverridecommandlockouts
% The preceding line is only needed to identify funding in the first footnote. If that is unneeded, please comment it out.
\usepackage{cite}
\usepackage{amsmath,amssymb,amsfonts}
\usepackage{verbatim}
\usepackage{optidef}
\usepackage{textgreek}
\usepackage{algorithmic}
\usepackage{graphicx}
\usepackage[export]{adjustbox}
\usepackage{textcomp}
\usepackage{xcolor}
%\newcommand\norm[1]{\left\lVert#1\right\rVert}
\newcommand\normx[1]{\left\Vert#1\right\Vert}
%\usepackage[margin=3cm]{geometry}
%\usepackage{algcompatible}
% OR \usepackage{algorithmic}
%\usepackage{algorithm}
\newcommand{\mybibliography}{\bibliography{jour_short,conf_short,IEEEexample.bib}}
\DeclarePairedDelimiterXPP\BigOSI[2]%
  {\mathcal{O}}{(}{)}{}%
  {\SI{#1}{#2}}
  
\def\BibTeX{{\rm B\kern-.05em{\sc i\kern-.025em b}\kern-.08em
    T\kern-.1667em\lower.7ex\hbox{E}\kern-.125emX}}
\begin{document}

\title{Deep Unfolding Enabled Constant Modulus Waveform Design for Joint Communications and Sensing\\
%{\footnotesize \textsuperscript{*}Note: Sub-titles are not captured in Xplore and should not be used}
%\thanks{P. Krishnananthalingam, N. T. Nguyen and M. Juntti are with Centre for Wireless Communications, University of Oulu, P.O.Box 4500, FI-90014, Finland (e-mails: \{prashanth.krishnananthalingam, nhan.nguyen, markku.juntti\}@oulu.fi).
}

\author{\IEEEauthorblockN{Prashanth Krishnananthalingam, Nhan Thanh Nguyen, Markku Juntti}%\IEEEauthorrefmark{1}\IEEEauthorrefmark{1}\IEEEauthorrefmark{1}
\IEEEauthorblockA{Centre for Wireless Communications, University of Oulu, P.O.Box 4500, FI-90014, Finland}%\IEEEauthorrefmark{1}
Emails: \{prashanth.krishnananthalingam, nhan.nguyen, markku.juntti\}@oulu.fi}
%
%
%
% \author{\IEEEauthorblockN{Prashanth Krishnananthalingam}
% \IEEEauthorblockA{\textit{CWC-Radio Technologies} \\
% \textit{University of Oulu}\\
% Oulu, Finland \\
% prashanth.krishnananthalingam@oulu.fi}
% \and
% \IEEEauthorblockN{Nhan Nguyen}
% \IEEEauthorblockA{\textit{CWC-Radio Technologies} \\
% \textit{University of Oulu}\\
% Oulu, Finland \\
% nhan.nguyen@oulu.fi}
% \and
% \IEEEauthorblockN{Markku Juntti}
% \IEEEauthorblockA{\textit{CWC-Radio Technologies} \\
% \textit{University of Oulu}\\
% Oulu, Finland \\
% markku.juntti@oulu.fi}
% }
%

\maketitle

%
\begin{abstract}
Joint communications and sensing (JCAS) systems have recently emerged as a promising technology to utilize the scarce spectrum in wireless networks and to reuse the same hardware to save infrastructure costs. In practical JCAS systems, dual functional constant-modulus waveforms can be employed to avoid signal distortion in nonlinear power amplifiers. However, the designs of such waveforms are very challenging due to the nonconvex constant-modulus constraint. The conventional branch-and-bound (BnB) method can achieve optimal solution but at the cost of exponential complexity and long run time. In this paper, we propose an efficient deep unfolding method for the constant-modulus waveform design in a multiuser multiple-input multiple-output (MIMO) JCAS system. The deep unfolding model has a sparsely-connected structure and is trained in an unsupervised fashion. It achieves good communications-sensing performance tradeoff while maintaining low computational complexity and low run time. Specifically, our numerical results show that the proposed deep unfolding scheme achieves a similar achievable rate compared to the conventional BnB method with 30 times faster execution time.
\end{abstract}

\begin{IEEEkeywords}
Joint communications and sensing, deep neural networks, deep unfolding, constant-modulus waveform design.
\end{IEEEkeywords}

\section{Introduction}
The developments in wireless communications systems create the need for more resources, especially for more spectrum which led to spectrum scarcity. This has driven the reuse of spectrum for both communications and radar sensing. The emerging concept of unifying communications and sensing can be referred to by different names, such as integrated sensing and communications (ISAC) \cite{ouyang2022performance, liu2022integrated} or joint communications and sensing (JCAS) \cite{zhang2018multibeam} which is the term used in this paper. %Joint communications and sensing (JCAS) systems enable the sharing of network resources between them by designing a single transmitted waveform for both purposes \cite{9585321}.
%Joint communications and sensing (JCAS) systems enable the sharing of network resources, by which a base station can communicate with mobile users and acts as a radar node simultaneously \cite{9411178}. %However, communications and sensing require different conditions to be satisfied in the waveform. For example, communications require high data rates while sensing requires a good correlation of the waveform \cite{9924202}. The ultimate goal of the design is to achieve the optimal waveform which should balance between the requirements of both communication and sensing \cite{9971740}.


Early JCAS designs focused on mitigating the mutual interference between the radar and communications subsystems to enable their smooth co-existence \cite{hassanien2016signaling, hassanien2015dual, liu2020joint, ma2020joint, chiriyath2017radar, zheng2019radar, li2016optimum, martone2019joint, liu2018mu, liu2017robust, buzzi2019using}. Specifically, Liu \textit{et al.} \cite{liu2018mu} investigated two multiple-input multiple-output (MIMO) JCAS systems based on sharing the transmit antennas among the communications data and radar probing signal transmission. Hassanien \textit{et al.} \cite{hassanien2015dual} optimized the dual functional waveform by transmitting communications data in the side lobes of the radar waveform, which is a radar-centric design and does not ensure a satisfactory communications performance in general. Wu \textit{et al.} \cite{wu2020waveform} proposed a frequency-hopping MIMO radar-based waveform for channel estimation.

Recent works focus on investigating MIMO JCAS systems to exploit spatial beamforming gains for a good communications--sensing performance tradeoff. The waveform designs were performed with various methods, such as manifold optimization with Riemannian conjugate gradient \cite{8288677}, and semi-definite relaxation \cite{9124713, 9940478}. These works focused on optimizing the transmit waveform under the communications performance constraint. Moreover, the radar Cramer-Rao lower bound \cite{9468975,9420261}, beam pattern matching error \cite{9104378}, and mutual information \cite{9145467} were also considered as the sensing metrics.  In \cite{liu2018toward, tang2022mimo, wu2023constant, liu2021dual}, the practical constant-modulus waveforms were particularly considered and optimized to avoid signal distortion in nonlinear power amplifiers in JCAS systems. Particularly, Liu \textit{et al.} \cite{liu2018toward} investigated the multiuser interference (MUI) minimization problem under the constant-modulus waveform constraint in multiuser MIMO JCAS systems. To tackle the challenging problem, the branch and bound (BnB) method has been employed to achieve optimal performance. However, its complexity is exponential with the number of transmit antennas and the run time is generally long. These challenges motivate us to leverage machine learning (ML), in particular deep learning (DL) techniques for constant-modulus waveform designs.

\begin{comment}
the computational cost of the algorithm was in exponential order with the number of transmit antennas and the number of users. In , the authors formulated the waveform design as an optimization problem under constant modulus power constraint and proposed a  method-based algorithm to solve for the optimal design. However, the computational cost of the algorithm is in exponential order of the number of transmit antennas. A  based algorithm was proposed in  to solve the waveform design problem. Even though the radar performance is improved, the design restricted the degrees of freedom (DoF) in the waveform design and ultimately affected the communication performance. Further, the authors noted that the radar beam pattern had a considerable distortion when the number of users is much smaller compared to the number of antennas.
\end{comment}

Machine learning has gained prominence in recent years in various fields of wireless communications \cite{8964474, 8642915, 9772315, 9047156, 9425589, 9626148}. Recently, DL has been leveraged for JCAS designs ~\cite{Mateos2022EndtoEnd,muth2023Autoencoder,xu2022deep,elbir2021terahertz}. Specifically, in \cite{Mateos2022EndtoEnd}, \cite{muth2023Autoencoder}, an end-to-end deep autoencoder model was proposed to learn the joint operations of the beamforming, target detection, and receiver processing. Whereas deep reinforcement learning-based JCAS designs were introduced by Xu \textit{et al.} \cite{xu2022deep} with the focus on sparse transmit arrays. Elbir {\it et al.} considered both model-based and data-based JCAS designs in \cite{elbir2021terahertz}. Specifically, convolutional neural networks and deep neural networks (DNNs) were trained to estimate the direction of the radar targets and for hybrid beamforming \cite{elbir2021terahertz}. Leveraging DL for JCAS enables reliable communications and sensing operations with fixed and low computational complexity and execution time. However, the existing DL-based JCAS designs mostly relied on the learning capabilities of black-box deep neural networks (DNNs). Consequently, they are not explainable as the conventional model-based optimization design methods. Furthermore, their training is generally a challenging task requiring much effort in fine tuning and generating massive data sets. Thus, many conventional optimizers and data-driven DNNs exhibit major limitations due to their resource-constraints, high complexity, and black-box nature. Alternatively, the deep unfolding approach leverage both domain knowledge and learning capabilities to build explainable DL models that achieve performance gains and are easier to implement \cite{9557819}. However, its application for JCAS has not yet been explored.

\begin{comment}
These algorithms are able to solve most statistical problems by accepting examples of inputs and their intended outputs. Developments in the field of machine learning led to a new sub-field called DL. DL algorithms are capable of solving the above-mentioned statistical problems via complicated algorithms widely known as deep architectures. They usually consist of neural networks which use many non-linear operations and layers. The applications are in diverse fields ranging from signal processing \cite{8678825, 9780487} and computer vision \cite{8243952,10074290,10.1007/978-981-19-0475-2_52} to climate science \cite{REN2021100178, ZHANG202053}. Neural Networks are applied in MIMO symbol detection as well. A DL network called DetNet was used for the MIMO detection problem \cite{8227772, 8642915}. Similar works which use deep neural networks for MIMO detection can be found in \cite{8780959,9772315,9047156}. The computational complexity of these neural networks is significantly reduced compared to the other conventional detection schemes.    
\end{comment}
%Several works deal with ML applications on JCAS systems in terms of beamforming \cite{9557817, 9492131,9791349}, channel estimation \cite{10001672}, radar sensing \cite{9827815}, receiver design \cite{9967989} and received signal demodulation \cite{ZHANG2023100574}. % However, there are no proper works on the application of neural networks for the purpose of JCAS waveform design to the best of our knowledge. On the other hand, 
 %These ML-based algorithms had low time and computational complexities compared to the conventional approaches while maintaining the same level of performance or outperforming them. % These results served as the motivation to focus on ML based waveform design for a JCAS system.
 %In \cite{9971740}, the authors  suggested applying artificial intelligence (AI) to waveform designs in JCAS systems. Deep unfolding \cite{9020494} is an ML-based algorithm, which is used extensively in many fields of communications including channel estimation \cite{9791132}, MIMO detection \cite{9860554}, and precoding \cite{8815519,9854863,9125984}. This method converts an iterative process with a fixed number of iterations into a deep neural network (DNN) by unfolding its structure and incorporating some trainable parameters. DL techniques are used in the learning process of these parameters. This DNN can be used for various problems \cite{9020494}. %Therefore, we consider applying a deep unfolding network for the purpose of JCAS waveform design. Applying a neural network-based approach to the waveform design problem could be challenging as a massive amount of data need to be processed. On the other hand, the prospects of complexity reduction act as a driving force to focus on this research direction.



In this paper, we consider a multiuser MIMO JCAS system. We aim at optimizing the communications-sensing performance tradeoff via minimizing the weighted sum of the MUI and the deviation from the benchmark sensing waveform. In particular, the constant-modulus waveform is considered, as motivated by the work of Liu \textit{et al.} \cite{liu2018toward}. To overcome the challenging design problem, we propose an efficient deep unfolding model based on the projected gradient descent method. The model is trained in an unsupervised strategy to output a reliable constant-modulus waveform. In particular, it has a sparsely connected network architecture, and, thus, it performs the JCAS design with very low computational complexity. %We overcome the high complexity  and derive an optimized transmit waveform that satisfies the communications and sensing requirements. We propose a weighted minimization problem between the MUI energy and the radar beam pattern similarity under the constant-modulus constraint. %The problem becomes non-convex due to the constant-modulus constraint and consequently NP-hard. First, we formulate the waveform design problem for a MIMO JCAS system as the trade-off minimization between the MUI energy and the radar beam pattern similarity under the constant modulus constraint (constant-modulus constraint). A strict similarity constraint on the radar beam pattern affects communication performance. Therefore, we allow a tolerable mismatch between the benchmark and designed radar beam patterns through this trade-off design. The problem becomes non-convex due to the constant-modulus constraint. 
%We develop a low-complexity deep unfolding DNN and train it with the channel matrix, desired symbol matrix, and benchmark radar signal matrix information. Then, we use the trained model to get the optimal transmitted symbol matrix. Further, we analyze the computational and time complexities of the proposed DNN. 
We finally present extensive simulation results of the proposed scheme in comparison with the conventional BnB approach \cite{8386661}. The simulation results demonstrate that the proposed unfolding network with a few layers and sparse neuron connections performs similarly to the BnB counterpart with significantly lower computational and time complexities.%  ies.

The remainder of this paper is organized as follows: The system model is described in section \ref{section:2}. Section \ref{section:3} proposes the deep unfolding waveform design and discusses the computational complexity. We provide the numerical results in section \ref{section:4} and conclude the paper in \ref{section:5}. 
\begin{figure}[htbp]
\centerline{\includegraphics[width=0.5\textwidth]{Figures/BS.png}}
\caption{Joint communications and sensing system.}
\label{fig:1}
\end{figure}
%\textit{Notations:} a, \textbf{a} and \textbf{A} denote a complex scalar value, a vector, and a matrix, respectively. \begin{math}[.]^T \end{math} and \begin{math}[.]^H \end{math} denote the transpose and conjugate transpose; \begin{math}\mathbb{C}\ \end{math} and \begin{math}\mathbb{R}\ \end{math} denote the sets of complex numbers and real numbers respectively. Moreover, \begin{math}\|.\|\end{math} and \begin{math}\|.\|_{F}\end{math} denote the standard \begin{math}\textit{l}_{2}\end{math} norm and Frobenius norm respectively.
\section{System Model}
\label{section:2}
We consider a multiuser MIMO JCAS system, where the base station (BS) equipped with $N$ antennas simultaneously transmits probing signals to the targets and data signals to $K$ %\textit{K} 
single-antenna users, as illustrated in Fig. 1. % Let $N$ denote the number of antennas at the BS.The channel matrix is denoted by \textbf{H} = [\textbf{h}_1 , \textbf{h}_2 , ..., \textbf{h}_K ]^T \in \begin{math}\mathbb{C}\ ^{{K}\times{N}}\end{math} with $\textbf{h}_j $ denotes the columns of \textbf{H}. The transmit signal matrix is denoted by \textbf{X} = [\textbf{x}_1, \textbf{x}_2,..., \textbf{x}_M ] \in \begin{math}\mathbb{C}\ ^{{N}\times{M}}\end{math}  where $ \textbf{x}_j $ denotes the columns of \textbf{X}, and $ M $ is the length of the communication frame. The noise matrix is denoted by \textbf{Z} = [\textbf{z}_1, \textbf{z}_2,..., \textbf{z}_M] \in \begin{math}\mathbb{C}\ ^{{K\times{M}}\end{math} where $\textbf{z}_j$ denotes the columns of \textbf{Z} and \textbf{z}_j \sim CN(0, N_0\textbf{I}_N ) , \forall j.
% \subsection{Communication Model}
%We use the following notations to denote the variables in the analysis.
Let $\vx_m \in \setC^{N \times 1}$ be the transmit signal vector in the $m$-th communications time frame, and let $\mX = [\vx_1, \ldots, \vx_M]$, with $M$ being the length of the communication frame. Take $\vh_n \in \setC^{K \times 1}$ as the channel vector corresponds to the $n$-th antenna of the BS, and consider $\mH = [\vh_1, \ldots, \vh_N]$. Consider $\vz_m \in \setC^{K \times 1}$ as the noise vector corresponds to the $m$-th communications time frame, where $\vz_m \sim CN(0, N_0\textbf{I}_N )$ and take $\mZ = [\vz_1, \ldots, \vz_M]$.

%\blue{You can write like: Let $\vx_m \in \setC^{N \times 1}$ be the transmit signal vector in the $m$-th communications time frame, and let $\mX = [\vx_1, \ldots, \vx_M]$... This doesn't mean that you apply the same sentence for all the notations, you must learn how to write....}

The received signal at $K$ users can be expressed as
\begin{equation}
\mY = \mH\mX + \mZ.\label{eq:1}
\end{equation}
\begin{comment}
where \textbf{H} = [\textbf{h}_1 , \textbf{h}_2 , ..., \textbf{h}_K ]^T \in \begin{math}\mathbb{C}\ ^{{K}\times{N}}\end{math} denotes the channel matrix with $\textbf{h}_j $ denotes the columns of \textbf{H},
\textbf{X} = [\textbf{x}_1, \textbf{x}_2,..., \textbf{x}_M ] \in \begin{math}\mathbb{C}\ ^{{N}\times{M}}\end{math} denotes the transmit signal matrix with $ \textbf{x}_j $ denotes the columns of \textbf{X}, and \textit{M} is the length of the communication frame,
\textbf{Z} = [\textbf{z}_1, \textbf{z}_2,..., \textbf{z}_M] \in \begin{math}\mathbb{C}\ ^{{K\times{M}}\end{math} is the noise matrix, with $\textbf{z}_j$ denotes the columns of \textbf{Z} and
\textbf{z}_j \sim CN(0, N_0\textbf{I}_N ) , \forall j.
\end{comment}
The transmit signal matrix $ \mX $ is used for both communications and radar sensing. Hence, each communications symbol is also a snapshot of a radar pulse. The downlink channel $ \mH $ is assumed to be flat Rayleigh fading and remains unchanged during a single communications frame/radar pulse. Moreover, the channel is assumed to be perfectly estimated at the BS via pilot symbols.

Let $\vs_m \in \setC^{K \times 1}$ be the desired symbol vector to be received by the users during the $m$-th communications time frame, and let $\textbf{S} = [\vs_1, \ldots, \vs_M]$, whose entries are assumed to be drawn from the same constellation for all users. 
 To unveil the effects of the MUI in the received signals, let us rewrite the signal model (\ref{eq:1}) as
\begin{equation}
\mY = \textbf{S} + (\mH\mX - \textbf{S})  + \mZ,\label{eq:2}
\end{equation}
where the second term, i.e., $(\mH\mX - \textbf{S})$, represents the MUI in the system. The total MUI power can be given as
\begin{equation}
P_{\textrm{MUI}} = \|\mH\mX - \textbf{S}\|_{F}^{2} = \sum_{i,j} |\vh_i^T\vx_j - s_{i,j}|^2, \label{eq:3}
\end{equation}
where $s_{i,j}$ is the (\emph{i, j})-th entry of \textbf{S}, and \begin{math}\mathbb{E}(.)\end{math} represents the ensemble average with respect to time index. The signal-to-interference-plus-noise ratio (SINR) per frame for the $i$-th user is given as
\begin{equation}
\gamma_{i} = \frac{\mathbb{E}(|s_{i,j}|^2)}{\mathbb{E}(|\vh_i^T\vx_j - s_{i,j}|^2) + N_0}. \label{eq:4}
\end{equation}
It is observed that the MUI directly impacts the SINR of the downlink users, which degrades the performance of the communications.

\subsection{Radar Model}%%%%%%%%%%%%%%%%%
%\red{You don't need this subsection.... it is unecessary and wrong...}

%Constant modulus waveforms are crucial for efficient transmission in both radar and communication applications, especially when low-cost non-linear power amplifiers are used \cite{5978417}. These waveforms prevent signal distortion and enable energy-efficient transmission. Researchers have extensively studied the use of constant modulus waveforms for massive MIMO communication \cite{6451071,7811286,7738555,8023970} scenarios and MIMO radar waveform designs \cite{4668417,6649991,7450660,7840054}. In these studies, optimization problems are formulated with a non-convex constant-modulus constraint. Therefore, we use a benchmark radar signal matrix $\textbf{X}_0$ with constant modulus entries.
We use the orthogonal chirp waveform matrix, which is adapted from \cite{7450660} as the benchmark waveform, $\mX_0$ for the radar system. Moreover, the transmit beam pattern of the JCAS system is given as%the transmit beam pattern of the JCAS system 
%For this radar benchmark signal matrix, first, the preliminary dual functional waveform matrix, \textbf{X}\textsubscript{P} associated with the spatial covariance matrix, \textbf{R}\textsubscript{X} is designed. The relationship between \textbf{X}\textsubscript{P} and \textbf{R}\textsubscript{X} is given as

%\begin{equation}
%\textbf{R}_X = \frac{1}{L}\ \textbf{X}_P\textbf{X}_P^H\label{eq:5}
%\end{equation}

%Then, element wise normalization is done to the resultant waveform matrix \textbf{X}\textsubscript{P} to obtain the desired waveform matrix \textbf{X}\textsubscript{0} with constant modulus entries.


\begin{equation}
P_\textrm{d}(\theta) = \frac{1}{M}\ \textbf{a}^{H}(\theta)\mX\mX^{H}\textbf{a}(\theta), \label{eq:5}
\end{equation}
where, $\theta$ refers to the detection angle and $ \textbf{a}(\theta) = [1, e^{j2\pi\Delta sin(\theta)}, \ldots, e^{j2\pi(N-1)\Delta sin(\theta)}] \in \setC^{N \times 1} $ is the antenna steering vector with $\Delta$ denoting the normalized antenna spacing with respect to the wavelength.

 %\section{Constant Modulus Waveform Design}
 %\label{section:3}
\subsection{Problem Formulation}
We consider a weighted minimization problem which can be formulated below.  Here, the objective function is the weighted sum of the MUI and the distance between $ \mX $ and the desired waveform. %The goal of this paper is to obtain a trade-off design which look after both communications and radar performances. The achievable sum rate is taken as the performance metric for the communications system, while the similarity between the desired and designed radar beam pattern is taken as the performance metric for radar system. The achievable sum rate can be maximized by minimizing the MUI energy. On the other hand, the designed radar beam pattern should match the desired radar beam pattern. Therefore, we formulate the optimization problem in (\ref{eq:7}) with weighted-sum of MUI and  between the MUI energy and the radar beam pattern similarity with constant modulus constraint.
\begin{mini!}%|s|[0]
    {\mX}{\rho\|\mH\mX -\textbf{S}\|_{F}^{2} + (1 - \rho)\|\mX - \mX_0\|_{F}^{2}  \label{eq:7a}}
    {\label{eq:7}}{}
  %  \addConstraint{\|vec(\textbf{X - X_{0}\)})\|_{F}^{2}}{\le \eta}
    \addConstraint{|x_{i,j}|}{ = \sqrt{\frac{P_{\textrm{T}}}{N}} \label{eq:7b}},
\end{mini!}
where $ P_{\textrm{T}} $ is the total transmit power, $\mX_0 \in \setC^{N \times 1} $ is the known benchmark radar signal matrix with constant modulus entries, $ \rho \in [0, 1] $ is the weighting factor for the communications systems, and $x_{i,j}$ is the (\emph{i, j})-th entry of $ \mX $. %There are two terms in the objective functions. The MUI term represents the optimization of communications while the similarity term corresponds to radar performance.
%The constant modulus constraint is used extensively in MIMO systems to improve the power gain by controlling only the phase of the symbols \cite{7523242}.

%\red{is this problem easy or challenging? and why is so? how it leads you to the reformulation and deep learning things?... Please provide the related discussion here.}
The feasible region of the optimization problem (\ref{eq:7}) is a circle of radius $ \sqrt{\frac{P_{\textrm{T}}}{N}} $ for each $ x_{i,j} $ due to the constant-modulus constraint. Thus, the problem is non-convex and NP-hard. %In \cite{6649991}, the authors solved a waveform optimization problem with constant-modulus constraint and reduced the complexity by using relaxation and randomization methods. However, the solution was approximate. Moreover, the similarity constraint also significantly impacts the sum rate \cite{9124713, 9827892}, as it forces a strict structure in the transmit waveform. \red{why you need this discussion? do the mentioned works considered the same problem as you? in that case, will you show that you get better performance than them by saying "to overcome the challenges..."?}
We next propose a deep unfolding waveform design based on a projected gradient descent approach.

\section{Proposed Deep Unfolding Waveform Design}
\label{section:3}
\subsection{Problem Reformulation}
We herein introduce the general idea to solve problem (\ref{eq:7}) based on the deep unfolding approach. To this end, we first rewrite the two terms in the objective function in (\ref{eq:7a}) as:
\begin{equation}
\|\mH\mX -\textbf{S}\|_{F}^{2} = \sum_{i=1}^M \|\mH\vx_i - \vs_i\|^{2}\label{eq:8}
\end{equation}
and
\begin{equation}
\|\mX - \mX_0\|_{F}^{2} = \sum_{i=1}^M \|\vx_i - \vx_{0i}\|^{2}, \label{eq:9}
\end{equation}
respectively. Here $ \vx_i $, $ \vs_i $ and $ \vx_{0i} $ are the $i$-th columns of $ \mX $, \textbf{S} and $ \mX_0 $ respectively. It is observed from (\ref{eq:8}) and (\ref{eq:9}) that the objective function is column-wise separable. Therefore, we consider an optimization problem which aims to optimize the columns of $ \mX $ separately for the corresponding columns of \textbf{S} and $ \mX_0 $, instead of (\ref{eq:7}). %\red{which problem?}
Furthermore, as the elements of $ \vx_i $ and $\vx_{0i}$ are constrained to have a constant amplitude, we consider normalized vector variables, $ \vx \in \setC^{N \times 1}$ and $ \vx_0 \in \setC^{N \times 1} $ such that $ \vx_i = \sqrt{\frac{P_{\textrm{T}}}{N}}\vx $ and $  \vx_{0i} = \sqrt{\frac{P_{\textrm{T}}}{N}}\vx_0 $. Here, we have also dropped the column index $ i $ for notational convenience. We reformulate the optimization problem (\ref{eq:7}) using these newly defined normalized vector variables as%Moreover, we convert the entries of \textbf{x} to be unit-modulus by dividing each entry of \textbf{x} by its amplitude and taking out \begin{math}\sqrt{\frac{P_{\textrm{T}}}{N}}\end{math}. , the optimization problemcan be reformulated 
\begin{mini!}
{\vx}{\rho \norm{\sqrt{\frac{P_{\textrm{T}}}{N}}\mH\vx-\vs}^{2} + (1 - \rho)\norm{\sqrt{\frac{P_{\textrm{T}}}{N}}(\vx - \vx_{0})}^{2} \label{eq:10a}}
{\label{eq:10}}{}
\addConstraint{|x(n)|}{=1 \label{eq:10b}}, \forall n,
\end{mini!}
where $ x(n) $ denotes the $n$-th entry of $\vx $. 

 Since problem (\ref{eq:7}) and (\ref{eq:10}) are equivalent in the sense that the latter is the column-wise reformulation of (\ref{eq:7}), we focus on (\ref{eq:10}) in the following analysis to develop the deep unfolding network. %\red{what is the "analysis of the neural network?"}

%\subsection{The structure of the neural network}
%Detnet is a neural network architecture that was proposed in \cite{8227772}. A detailed analysis for various modulation schemes and comparisons with other algorithms are provided in\cite{8642915}. It is a multi-layer neural network dedicated to MIMO detection.
\begin{comment}
We utilize the ScNet proposed in \cite{8780959} for massive MIMO detection in our waveform design problem. ScNet was the improved version of DetNet discussed in \cite{8227772}. We, first describe the MIMO detection problem to clarify how the ScNet can be incorporated into the waveform design problem. Let's consider the standard linear MIMO model with an E-antenna base station and F users as described by equation (\ref{eq:11}).

\begin{equation}
    \boldsymbol{\bar{y}} = \boldsymbol{\bar{H}}\boldsymbol{\bar{x}} + \boldsymbol{\bar{w}}\label{eq:11}
\end{equation}

where \boldsymbol{\bar{y}} \in \begin{math}\mathbb{C}\ ^{{E}}\end{math} is the received vector, \boldsymbol{\bar{H}} \in \begin{math}\mathbb{C}\ ^{{E}\times{F}}\end{math} is the channel matrix, \boldsymbol{\bar{x}} \in \begin{math}\mathbb{S}\ ^{{F}}\end{math} is the unknown vector from a finite constellation \begin{math}\mathbb{S}\ \end{math}, and \boldsymbol{\bar{w}} \in \begin{math}\mathbb{C}\ ^{{E}}\end{math} is the noise vector consists of elements from a zero mean complex random normal variable.

The estimation of the transmitted symbols vector through the maximum likelihood (ML) detection can be realized as a minimum distance problem described in equation (\ref{eq:12}).

\begin{mini}|s|[0]
    {\boldsymbol{\bar{x}}}{\|\boldsymbol{\bar{y}} -\boldsymbol{\bar{H}}\boldsymbol{\bar{x}}\|^{2}}
    {}
    {\label{eq:12}}{}
\end{mini}

However, these variables cannot be used in neural networks directly as neural networks generally work with real numbers. Therefore, the above variables must be converted to equivalent real-valued variables. The equivalent real-valued problem can be in the form of equation (\ref{eq:13})

\begin{mini}|s|[0]
    {\textbf{x}}{\|\textbf{y} - \textbf{H}\textbf{x}\|^{2}}
    {}
    {\label{eq:13}}{}
\end{mini}
where, \textbf{y} = 
  \begin{bmatrix}
    \Re(\boldsymbol{\bar{y}})\\ 
    \Im(\boldsymbol{\bar{y}})
  \end{bmatrix} \in \begin{math}\mathbb{R}\ ^{{2E}}\end{math}, \textbf{x} = 
  \begin{bmatrix}
    \Re(\boldsymbol{\bar{x}})\\ 
    \Im(\boldsymbol{\bar{x}})
  \end{bmatrix} \in \begin{math}\mathbb{R}\ ^{{2F}}\end{math}, and \textbf{H} = 
  \begin{bmatrix}
    \Re(\boldsymbol{\bar{H}}) & -\Im(\boldsymbol{\bar{H}})\\ 
    \Im(\boldsymbol{\bar{H}}) & \Re(\boldsymbol{\bar{H}})
  \end{bmatrix} \in \begin{math}\mathbb{R}\ ^{{2E}\times{2F}}\end{math}
  
The modified problem (\ref{eq:13}) was taken for the estimation through DetNet. DetNet is a fully-connected network and the connection diagram of the DetNet for a 2-by-2 MIMO system after the expansion is depicted in Figure 2. There are four inputs in each layer, namely, $\textbf{H}^T$\textbf{y}, $\textbf{H}^T$\textbf{H}, $\textbf{x}_{in}^{l}$ and $\textbf{v}_{in}^{l}$. There are two outputs $\textbf{x}_{out}^{l}$, the estimation of the unknown transmitted vector, and $\textbf{v}_{out}^{l}$ where \textit{l} represents the layer index.
 
\begin{figure}[htbp]
% \centerline{\includegraphics[scale=0.5]{BS.png}}%[max size={\textwidth}{\textheight}]
\centerline{\includegraphics[width=0.5\textwidth]{DetNet.png}}%[max size={\textwidth}{\textheight}]
\caption{The network connection diagram of the DetNet.}
\label{fig:2}
\end{figure}

The authors of \cite{8780959}, analyzed the network structure of the DetNet and simplified the inputs and outputs by removing the input vector $\textbf{v}_{in}^{l}$ and the output vector $\textbf{v}_{out}^{l}$ as they did not carry any information. Similarly, they modified the fully-connected structure by connecting only the same indexed vector elements to the output. Figure 3 shows the connection diagram of the ScNet, which resulted from these modifications.

\begin{figure}[htbp]
% \centerline{\includegraphics[scale=0.5]{BS.png}}%[max size={\textwidth}{\textheight}]
\centerline{\includegraphics[width=0.5\textwidth]{ScNet.png}}%[max size={\textwidth}{\textheight}]
\caption{The network connection diagram of the ScNet.}
\label{fig:3}
\end{figure}

We compare the equation (\ref{eq:12}) and the first part of the equation (\ref{eq:10a}) and realize that these two equations are mathematically similar. The standard ML detection problem tries to estimate \textbf{x} such that the Euclidean distance between the \textbf{y} and \textbf{Hx} becomes minimal. Similarly,
\end{comment}
\begin{comment}The operation of the neural network can be described as a minimization process of variables based on some objective function. Similarly, the waveform design problem is a minimization problem as explained in equation (\ref{eq:10}). However, the variables in (\ref{eq:10}) cannot be used in neural networks directly as neural networks generally work with real numbers. Therefore, the above variables must be converted to equivalent real-valued variables.  The waveform optimization problem can be thought of as estimating \textbf{x} which results in the minimization of MUI in the system. The variable \textbf{y} in the neural network corresponds to \textbf{s} for the waveform design problem. \end{comment}
Because neural networks generally work with real-valued data, we present the equivalent real-valued problem of (\ref{eq:10}) as follows:
\begin{mini}
{\boldsymbol{\bar{x}}}{\rho\norm{\sqrt{\frac{P_{\textrm{T}}}{N}}\boldsymbol{\bar{H}\bar{x}}-\boldsymbol{\bar{s}}}^{2} + (1 - \rho)\norm{\sqrt{\frac{P_{\textrm{T}}}{N}}(\boldsymbol{\bar{x}} - \boldsymbol{\bar{x}}_{0})}^{2},}
{}{} \label{eq:11}
\end{mini}
\begin{comment}
\begin{mini!}
{\textbf{x}}{\rho\|\sqrt{\frac{P_{\textrm{T}}}{N}}\textbf{Hx}-\textbf{s}\|^{2} + (1 - \rho)\|\sqrt{\frac{P_{\textrm{T}}}{N}}(\textbf{x} - \textbf{x}_{0})\|^{2} \label{eq:10a}}
{\label{eq:10}}{}
\addConstraint{|x(n)|}{=1 \label{eq:10b}}, \forall n,
\end{mini!}

\begin{mini}|s|[0]
    {\boldsymbol{x}}}{\rho\|\sqrt{\frac{P_T}{N}}\boldsymbol{\bar{H}}\boldsymbol{\bar{x}}-\boldsymbol{\bar{s}}\|^{2} + (1 - \rho)\|\sqrt{\frac{P_T}{N}}(\boldsymbol{\bar{x}} - \boldsymbol{\bar{x}}_{0})\|^{2}}
    {}
    {\label{eq:11}}{}
\end{mini}
\end{comment}
where, $ \boldsymbol{\bar{s}} = 
  \begin{bmatrix}
    \Re(\vs)\\ 
    \Im(\vs)
  \end{bmatrix} \in \mathbb{R}\ ^{{2K}},
\boldsymbol{\bar{x}} = 
  \begin{bmatrix}
    \Re(\vx)\\ 
    \Im(\vx)
  \end{bmatrix} \in \mathbb{R}\ ^{{2N}},
\boldsymbol{\bar{x}}_{0} = 
  \begin{bmatrix}
    \Re(\vx_{0})\\ 
    \Im(\vx_{0})
  \end{bmatrix} \in \mathbb{R}\ ^{{2N}},  and
\boldsymbol{\bar{H}} = 
  \begin{bmatrix}
    \Re(\mH) & -\Im(\mH)\\ 
    \Im(\mH) & \Re(\mH)
  \end{bmatrix} \in \mathbb{R}\ ^{{2K}\times{2N}} $
\begin{comment}
Neural networks estimate the output based on an objective function that must be minimized. It is called the loss function. We need to minimize the value of the objective function (\ref{eq:11}) for our problem. Therefore, we supply  equation (\ref{eq:11}) as the loss function to the neural network and receive the optimal $ \textbf{x}_{opt} $ as the output from neural network. Further, we ensure the unit modulus constraint given in (\ref{eq:10b}) by calculating the absolute values of the elements of complex \textbf{x} and dividing the corresponding elements of resultant real \boldsymbol{\hat{x}} from the output by it in each layer before feeding it to the next layer. The optimization process in the neural network is based on a projected gradient-descent method. The iterative procedure of this method estimate the new value from the previous iteration based on equation (\ref{eq:12}).
\end{comment}

In light of the projected gradient-descent method, the iterative update procedure of the waveform in the $p$-th iteration can be expressed as
\begin{equation}
\boldsymbol{\hat{x}}_{p} =\prod \left(\boldsymbol{\hat{x}}_{p-1} - \delta_{p-1} \frac{\partial f(\boldsymbol{\bar{x}})}{\partial \boldsymbol{\bar{x}}} \vline_{\boldsymbol{\bar{x}} = \boldsymbol{\hat{x}}_{p-1}}\right), \label{eq:12} %; \boldsymbol{\hat{x}} (\boldsymbol{\bar{H}},\boldsymbol{\bar{s}})
\end{equation}   
%\red{you used $K$ to denote the number of user already, so $k$ usually mean the $k$-th user....}
%\red{...and don't use bold symbols for the index....}
\begin{comment}
\begin{equation}
\boldsymbol{\hat{x}_{k+1}} =\prod (\boldsymbol{\hat{x}_{k}} - \delta_{k} \frac{\partial \|\sqrt{\frac{P_T}{N}}\boldsymbol{\bar{H}}\boldsymbol{\bar{x}}-\boldsymbol{\bar{s}}\|^{2}}{\partial \boldsymbol{\bar{x}}} \vline_{\boldsymbol{\bar{x}} = \boldsymbol{\hat{x}_{k}}}) \label{eq:12}
\end{equation}
\end{comment}
where $\boldsymbol{\hat{x}}_{p}$ denotes the output vector of the neural network in the $p$-th iteration, $ \delta_{p-1} $ denotes the step size, and $ \Pi (\cdot) $ denotes the nonlinear projection operator and  $f(\cdot)$ is the objective function in (\ref{eq:11}), i.e., $ \rho\norm{\sqrt{\frac{P_{\textrm{T}}}{N}}\boldsymbol{\bar{H}\bar{x}}-\boldsymbol{\bar{s}}}^{2} + (1 - \rho)\norm{\sqrt{\frac{P_{\textrm{T}}}{N}}(\boldsymbol{\bar{x}} - \boldsymbol{\bar{x}}_{0})}^{2} $. %\red{use $\cdot$, not period. use symbols dollar for math equation, they are not text}.
The iterative update procedure in (\ref{eq:12}) motivates the application of a DNN of $ L $ layers to output the transmit waveform. Specifically, we aim at developing a deep unfolding DNN architecture mimicking the nonlinear transformation in (\ref{eq:12}). As long as the DNN is well constructed and trained, it can generate the solution vector $ \boldsymbol{\bar{x}} $ of the problem (\ref{eq:11}) for the given weight factor $ \rho $. %\red{so the communications performance can be poor?} 
We detail the structure and training of the DNN next.
%
\subsection{Deep Unfolding Model's Structure and Training}
From (\ref{eq:11}) and (\ref{eq:12}) we have
\begin{multline}
\boldsymbol{\hat{x}}_{p} =\prod \Bigg( \Bigg.\boldsymbol{\hat{x}}_{p-1} - \delta_{p-1} \Bigg( \Bigg.2\rho \frac{P_{\textrm{T}}}{N} \boldsymbol{\bar{H}}^T\boldsymbol{\bar{H}}\boldsymbol{\hat{x}}_{p-1} - 2\rho \sqrt{\frac{P_{\textrm{T}}}{N}} \boldsymbol{\bar{H}}^T\boldsymbol{\bar{s}}\\ + 2(1-\rho) \frac{P_{\textrm{T}}}{N} \boldsymbol{\hat{x}}_{p-1} - 2(1-\rho) \frac{P_{\textrm{T}}}{N} \boldsymbol{\bar{x}}_{\textrm{0}}\Bigg. \Bigg) \Bigg. \Bigg). \label{eq:14}
\end{multline}
\begin{comment}
\begin{equation}
%\|\textbf{X} - \textbf{X}_0\|_{F}^{2} = \sum_{i=1}^M \|\textbf{x}_i - \textbf{x}_{0i}\|^{2}\label{eq:9}
l( log(l)(\rho)
\end{equation}
\end{comment}
%
%The training process of neural network is difficult due to various reasons such as 
%
%Hence, the iterative update procedure given in (\ref{eq:12}}) can be re-written as,
Here, $\boldsymbol{\hat{x}}_{p}$ is the estimated transmit signal vector in the $p$-th iteration. Equation (\ref{eq:14}) shows that it depends on $\boldsymbol{\bar{x}_{0}}$, $\boldsymbol{\bar{H}}^T\boldsymbol{\bar{s}}$, $\boldsymbol{\bar{H}}^T\boldsymbol{\bar{H}}$ and $\boldsymbol{\hat{x}}_{p-1}$. Therefore, these variables are used as inputs to the deep unfolding network. Fig. \ref{fig:2} shows the architecture of the proposed neural network. There are four vector inputs in each layer, namely, $\boldsymbol{\bar{x}_{0}}$, $\boldsymbol{\bar{H}}^T\boldsymbol{\bar{s}}$, $\boldsymbol{\bar{H}}^T\boldsymbol{\bar{H}}$ and $\boldsymbol{\hat{x}}_{p-1}$ %\red{all are just bold-face text, not math symbols, and different from their defnitions.}
and the vector output of the corresponding layer is $\boldsymbol{\hat{x}}_{p}$, the estimation of the unknown transmitted vector. The weight matrices for each input are denoted by $\boldsymbol{W}_{1}^{p}$, $\boldsymbol{W}_{2}^{p}$, $\boldsymbol{W}_{3}^{p}$ and $\boldsymbol{W}_{4}^{p}$. The bias vectors for each input are denoted by $\boldsymbol{b}_{1}^{p}$, $\boldsymbol{b}_{2}^{p}$, $\boldsymbol{b}_{3}^{p}$ and $\boldsymbol{b}_{4}^{p}$. By inspecting (\ref{eq:14}), only the same indexed elements of the input and output vectors are related. Therefore, the neural network has a sparsely connected structure, where the input and the output are connected element-wise in each layer. This helps to reduce the computational complexity significantly. Due to the sparse connections, the weight matrices become diagonal matrices. The activation function is  denoted by $\psi(t)$, which is defined as
\begin{equation}
\psi(t) = -1 + 2(\sigma(t+0.5)-\sigma(t-0.5)),
\end{equation}
where $\sigma(t)$ denotes the rectified linear unit (ReLU) function.
The output of the $p$-th layer is updated as follows:

\begin{equation}
\boldsymbol{m}_{1}^{p} = \boldsymbol{W}_{1}^{p}\boldsymbol{\bar{x}}_{0} + \boldsymbol{b}_{1}^{p}, \label{eq:15}
\end{equation}
\begin{equation}
\boldsymbol{m}_{2}^{p} = \boldsymbol{W}_{2}^{p}\boldsymbol{\bar{H}}^T\boldsymbol{\bar{s}} + \boldsymbol{b}_{2}^{p}, \label{eq:16}
\end{equation}
\begin{equation}
\boldsymbol{m}_{3}^{p} = \boldsymbol{W}_{3}^{p}\boldsymbol{\bar{H}}^T\boldsymbol{\bar{H}\hat{x}}_{p-1} + \boldsymbol{b}_{3}^{p}, \label{eq:17}
\end{equation}
\begin{equation}
\boldsymbol{m}_{4}^{p} = \boldsymbol{W}_{4}^{p}\boldsymbol{\hat{x}}_{p-1} + \boldsymbol{b}_{4}^{p}, \label{eq:18}
\end{equation}
\begin{equation}
\boldsymbol{s}_p = \boldsymbol{m}_{1}^{p} + \boldsymbol{m}_{2}^{p} + \boldsymbol{m}_{3}^{p} + \boldsymbol{m}_{4}^{p}, \label{eq:19}
\end{equation}
\begin{equation}
\boldsymbol{\hat{x}}_p = \psi(\boldsymbol{s}_p). \label{eq:20}
\end{equation}

\begin{figure}[htbp]
\centerline{\includegraphics[width=0.5\textwidth]{Figures/Complete_NN.png}}
\caption{Architecture of the neural network.}
\label{fig:2}
\end{figure}

To avoid the vanishing gradients, the sensitivity to initialization, and the saturation of activation functions, we leverage a loss function that covers the output from all the layers. This technique is inspired by the idea of auxiliary classifiers discussed in GoogLeNet \cite{7298594}. The loss function can be given as,% in equation (\ref{eq:13}).
\begin{multline}
  %\|\textbf{X} - \textbf{X}_0\|_{F}^{2} = \sum_{i=1}^M \|\textbf{x}_i - \textbf{x}_{0i}\|^{2}\label{eq:9}
%\resizebox{.9\hsize}{!}{
l( \boldsymbol{x}_{opt} ; \boldsymbol{\hat{x}} \boldsymbol{(\bar{H},\bar{s})}) = \sum_{p=1}^L \Bigg( \Bigg.\rho\norm{\sqrt{\frac{P_T}{N}}\boldsymbol{\bar{H}}\boldsymbol{\hat{x}}_{p}-\boldsymbol{\bar{s}}}^{2} \\ + (1 - \rho)\norm{\sqrt{\frac{P_T}{N}}(\boldsymbol{\hat{x}}_{p} - \boldsymbol{\bar{x}}_{0})}^{2}\Bigg. \Bigg).\label{eq:13}  
\end{multline}
\begin{comment}
We train the neural network to obtain each column of \textbf{X} separately for a particular channel realization and the corresponding desired symbol vector. These columns are concatenated properly to obtain the final transmitted symbol matrix \textbf{X}. However, the neural network is capable of optimizing the entire batch of channel realizations at once. This speeds up the estimation process and it is also evident from the table \ref{tab:1}.% and \ref{tab:2}.
%
\section{Complexity Analysis}
This section is dedicated to the complexity analysis of the BnB algorithm, which was proposed in \cite{8386661} for the constant modulus waveform design, and the neural network architecture. First, the computational complexity is taken into the consideration. Then, the time complexity of both algorithms is investigated.   
\end{comment}
\subsection{Computational Complexity}
\label{subsec:3a}
We first note that $\boldsymbol{\bar{H}}^T\boldsymbol{\bar{s}}$ and $\boldsymbol{\bar{H}}^T\boldsymbol{\bar{H}}$ are computed beforehand and fed to the network.
\begin{comment}
\begin{figure}[htbp]
% \centerline{\includegraphics[scale=0.5]{BS.png}}%[max size={\textwidth}{\textheight}]
\centerline{\includegraphics[width=0.5\textwidth]{Netstr.png}}}%[max size={\textwidth}{\textheight}]
\caption{The network structure of the neural network.}
\label{fig:3}
\end{figure}
\end{comment}
Since $\boldsymbol{\bar{H}}^T\boldsymbol{\bar{H}} \in $ \begin{math}\mathbb{R}\ ^{{2N}\times{2N}}\end{math}, computing $ \boldsymbol{\bar{H}}^T\boldsymbol{\bar{H}}\boldsymbol{\hat{x}}_{p-1} $ requires a total of $ 2N(4N-1) $ floating point operations (FLOPs). The multiplications in (\ref{eq:15}) - (\ref{eq:18})   need  $ 8N $ FLOPs. Furthermore, additions in (\ref{eq:15}) - (\ref{eq:18})   need  $ 8N $ FLOPs. The additions in (\ref{eq:19}) need $ 6N $ FLOPs. The total arithmetic operations per layer will be $ 2N(4N+11) $ FLOPs. Hence, the computational complexity is $\mathcal{O} (2N(4N+11)L)$. As the computational complexity of the BnB is exponential with $ N $, it is higher compared to the proposed deep unfolding scheme. \cite{Tuy2016}

\section{Numerical Results}
\label{section:4}
In this section, we present the numerical results to validate the proposed deep unfolding network. %neural network \red{again...neural network...}. 
The total power budget is set to $ P_{\textrm{T}} = 30$ $ dBm $ and the channel is assumed to be flat Rayleigh fading. Moreover, the losses at the transmitter and receiver are assumed to be cancelled by the amplifier gains at the transmitter and receiver respectively. %\red{I suggested you not to use italic for acronym...}. 
Each entry of the channel matrix $\mH$ is assumed to follow the standard complex Gaussian distribution. %\red{this is your assumption, you write like it always have that distribution}
Moreover, we set  the number of transmit antennas as $ N $ = 8 and assume that the ULA with half-wavelength spacing is deployed at the BS. The desired symbol matrix \textbf{S} is assumed to be populated with symbols from the unit-power QPSK alphabet. %\red{this is assumption}
The length of the communication frame/radar pulse is set as $ M $ = 20. The number of layers for the neural network is set as $L$ = 10. %\red{again, don't use textit, use dollars}
We train the network with batches of 100 channel realizations. For  the benchmark radar signal, a constant modulus chirp signal is used. The noise signal is assumed to be complex additive white Gaussian (AWG) with zero mean and the variance of $N_0$. The transmit SNR is defined as SNR = \begin{math}P_{\textrm{T}}/N_0\end{math}.

The neural network is implemented in Python with the Tensorflow library. An initial learning rate of 0.0001 and a decaying factor of 0.97 are used. Adam optimizer is used to optimize the neural network. The input $\boldsymbol{\hat{x}}_{\textrm{0}}$ is initialized as a vector of zeros.

\begin{figure}[htbp]
% \centerline{\includegraphics[scale=0.5]{BS.png}}%[max size={\textwidth}{\textheight}]
\centerline{\includegraphics[width=0.5\textwidth]{Figures/AASRvsSNR.pdf}}%[max size={\textwidth}{\textheight}]
\caption{Average achievable sum rate with SNR for different approaches.}
\label{fig:3}
\end{figure}

\begin{comment}
First, we train our neural network for different \(\rho\) values and get the neural network model for that particular \(\rho\). Then, we test the trained models with a batch of test data. After that, we receive the batch of optimized transmitted symbol matrix \textbf{X} as the output from the neural network model. Finally, we evaluate the performance metrics using the output data set from the neural network model. We investigate the performance of the communications system by observing the sum rate response. 
\end{comment}

Fig. \ref{fig:3} shows the system average achievable sum rate with SNR for different \(\rho\) and \(\epsilon\). Here, \(\epsilon\) corresponds to normalized tolerable difference defined in \cite{8386661}. %We observe the radar beam pattern to justify the performance of the radar system. Fig. \ref{fig:4} shows the radar beam pattern with respect to the reference radar beam pattern. By observing the results in Fig. \ref{fig:3} and \ref{fig:4}, it can be stated that the proposed neural network provides a balanced trade-off design by considering both the sum rate and radar waveform similarity into the consideration. The sum rate gets improved with $ \rho $, while the radar beam pattern gets degraded.
%It can be observed from Fig. \ref{fig:3} that the sum rate improves quickly when \(\rho\) is higher. 

\begin{figure}[htbp]
% \centerline{\includegraphics[scale=0.5]{BS.png}}%[max size={\textwidth}{\textheight}]
\centerline{\includegraphics[width=0.5\textwidth]{Figures/beampattern.pdf}}%[max size={\textwidth}{\textheight}]
\caption{Radar beam patterns obtained by different approaches.}
\label{fig:4}
\end{figure}

Fig. \ref{fig:4} shows the system average radar beam pattern for different \(\rho\) and \(\epsilon\). The reference beam pattern is also included for the comparison. %We observe the radar beam pattern to justify the performance of the radar system. Fig. \ref{fig:4} shows the radar beam pattern with respect to the reference radar beam pattern. By observing the results in Fig. \ref{fig:3} and \ref{fig:4}, it can be stated that the proposed neural network provides a balanced trade-off design by considering both the sum rate and radar waveform similarity into the consideration. The sum rate gets improved with $ \rho $, while the radar beam pattern gets degraded.
It can be observed from Fig. \ref{fig:4} that the designed beam pattern through the neural network tends to match with the reference waveform when \(\rho\) is lower. Moreover, by comparing the results of neural network and BnB, it can be stated that the neural network keeps the main lobe in the desired direction even though the transmit power gets reduced while the beam pattern gets a significant degradation for BnB. Furthermore, by observing the sum rate and beam pattern results for \(\rho\) = 0.8 and \(\epsilon\) = 1.5, we could state that the neural network maintain both a better beam pattern and better sum rate response compared to BnB.

We use the results for the constant modulus waveform design in \cite{8386661} through the BnB algorithm as a benchmark to validate our results. %However, the trade-off constants are defined differently from our analysis. Therefore, we devise a mechanism for the comparison. We consider the average mean squared error (MSE) between the benchmark and designed beam pattern as the performance metric for the radar system. On the other hand, we choose the average achievable sum rate as the performance metric for the communications system. Then, we obtain these two metrics from both algorithms for different \(\rho\) and  \(\epsilon\) values. Finally, we compare the variation of sum rate with average MSE for both algorithms. Fig. \ref{fig:5} shows the variation of sum rate with average MSE of the radar beam pattern for different SNR values.
We compare the system average sum rate with the average radar mean squared error (MSE) for both methods. Fig. \ref{fig:5} shows the variation of sum rate with average MSE of the radar beam pattern for different SNR values. The variation in Fig. \ref{fig:5} implies that the proposed neural network performs almost the same as the BnB algorithm. Moreover, a closer inspection of the results shows that the results of neural network are slightly better in most cases.
\begin{figure}[htbp]
% \centerline{\includegraphics[scale=0.5]{BS.png}}%[max size={\textwidth}{\textheight}]
\centerline{\includegraphics[width=0.5\textwidth]{Figures/AASRvsMSE.pdf}}%[max size={\textwidth}{\textheight}]
\caption{Variation of average achievable sum rate with average MSE of the radar beam pattern.}
\label{fig:5}
\end{figure}

We analyze the time complexities of both methods next. We compare their run times in the same machine. The BnB algorithm was implemented and tested through MATLAB. We kept all the values of the variables except $N$ as same as the values mentioned at the beginning of the section. Table \ref{tab:1} shows the average run time per channel for both methods with different transmit antennas. We consider the communication-only case for this purpose. The per-channel execution times for the neural network are significantly lower compared to the BnB. Based on the results, it can be stated that the time complexity of the neural network is lower than BnB.

\begin{comment}

%\subsection{Time complexity}
Further, the time complexity of the neural network is compared with the BnB algorithm as well. The execution times for the BnB algorithm and the neural network are used for this purpose. Here, the execution time for the neural network refers to the time taken for the trained model to complete the testing process. All the measurements were obtained by running the algorithms in the same server which uses 4 Intel(R) Xeon(R) Gold 6244CPU @ 3.60GHz Processors. The operating system was 64-bit Windows Server 2019 Standard and it had a memory of 384 GB RAM. Table \ref{tab:1} shows the comparison of execution times of both methods for different batch sizes.

%\begin{comment}
We set N = 8, K = 4, M = 20, and L = 10 
during all the simulations 

for the comparison of execution times.


Table \ref{tab:1} shows the execution times per channel of both algorithms for different batch sizes of channel realizations under the communication-only case where  \(\rho\) = 1 and  \(\epsilon\) = 2.  Here, \(\epsilon\) refers to the normalized tolerable difference mentioned in section V of \cite{8386661}. %Similarly, table \ref{tab:2} depicts the execution times under the sensing-only case where  \(\rho\) = 0 and  \(\epsilon\) = 0.
\end{comment}
\begin{table}[htbp]
\caption{Run times for both BnB and NN methods with number of transmit antennas}
\begin{center}
\begin{tabular}{|c|c|c|}
\hline
\textbf{Transmit}&\multicolumn{2}{|c|}{\textbf{Algorithm}} \\
\cline{2-3} 
\textbf{Antennas} & \textbf{\textit{Branch and Bound}}& \textbf{\textit{Neural Network}} \\
\hline
8 & 15.9 s & 0.507 s  \\
\hline
16 & 17.4 s & 0.532 s  \\
\hline
\end{tabular}
\label{tab:1}
\end{center}
\end{table}

\begin{comment}
\begin{table}[htbp]
\caption{Execution time with batch size for the algorithms under sensing-only case}
\begin{center}
\begin{tabular}{|c|c|c|}
\hline
\textbf{Batch}&\multicolumn{2}{|c|}{\textbf{Algorithm}} \\
\cline{2-3} 
\textbf{Size} & \textbf{\textit{Branch and Bound}}& \textbf{\textit{Neural Network}} \\
\hline
10 & 4.92 s & 122 s  \\
\hline
20 & 34.1 s & 94.0 s  \\
\hline
50 & 19.8 s & 182 s  \\
\hline
100 & 36.6 s & 99 s  \\
\hline
\end{tabular}
\label{tab:2}
\end{center}
\end{table}

The results from table \ref{tab:1} \begin{comment}
 and \ref{tab:2}  
 clearly shows that the neural network results converged quicker than the BnB algorithm. Moreover, the per-channel execution time of the neural network drops significantly with the batch size, while the execution time for the BnB algorithm remains almost the same. This trend shows that our neural network performs much better when the batch size of channel realizations increased. This can be stated as the distinct power of the neural network compared to the BnB algorithm, where the entire batch can be trained at once.

We also noted the convergence times for the sensing-only case where  \(\rho\) = 0 and  \(\epsilon\) = 0. The BnB algorithm converged quicker than the neural network in this case. However, we emphasize the point that, for the sensing-only case, the feasible region for the BnB-based optimization problem becomes a single point in the Euclidean space. Hence, the direct solution for this problem would be $\textbf{x}_0$ itself. On the other hand, the neural network will try out all the possible vectors with unit-modulus elements before converging. Therefore, the neural network needs a considerably higher time than the BnB. However, waveform design is not needed for the sensing-only case as the transmitted signal matrix \textbf{X} becomes the benchmark radar signal matrix, $\textbf{X}_0$.
\end{comment}
%To summarize the results,%
%. Moreover, the analysis of complexity of our neural network indicates that there is a significant complexity reduction compared to the BnB. Therefore, our proposed neural network for the purpose of waveform design in JCAS systems performs at the same level as the BnB algorithm complexity reductions in terms of both computational cost and execution time

\section{Conclusions}
\label{section:5}
In this paper, we proposed a neural network-based solution to the problem of waveform optimization in JCAS systems. Considering the results in section \ref{section:4} along with the complexity analysis in section \ref{subsec:3a}, it can be stated that our scheme gives almost the same performance as the BnB algorithm while experiencing reductions in both computational and time complexities. %The neural network scheme experiences a significant reduction of time complexity with respect to BnB when waveform optimization needs to be done for a large batch of channel realizations.
Simulation results validated the claim that our neural network performs as same as the BnB.%and time complexities are 

\section*{Acknowledgement}
This research has been supported by the Academy of Finland, 6G Flagship program and Infotech Oulu under Grant 346208.

\bibliographystyle{IEEEtran}
\mybibliography

\vspace{12pt}

\end{document}
